\subsection{埃尔米特型}

我们回顾《代数》中给出的下述定义（A, IX, § 3, No.~1）：

定义 1. --- 设 E 为域 K 上的向量空间。E 上的（左）埃尔米特型是从 E × E 到 K 中的一个映射 f，它满足下述条件（对 E 中的 \(x_{1}\) 、 \(x_{2}\) 、x、\(y_{1}\) 、 \(y_{2}\) 、y 与 K 中的 \(\lambda\) 、 \(\mu\) ）：

\begin{enumerate}
\def\labelenumi{(\arabic{enumi})}
\tightlist
\item
\item
\end{enumerate}

\[
\begin{array}{c} \left\{ \begin{array}{l} f (x _ {1} + x _ {2}, y) = f (x _ {1}, y) + f (x _ {2}, y) \\ f (x, y _ {1} + y _ {2}) = f (x, y _ {1}) + f (x, y _ {2}) \end{array} \right. \\ \left\{ \begin{array}{l} f (\lambda x, y) = \overline {{\lambda}} f (x, y) \\ f (x, \mu y) = \mu f (x, y) \\ f (x, y) = \overline {{f (y , x)}}  . \end{array} \right. \end{array}\tag{3}
\]

当域 K 是 R 时，E 上的埃尔米特型的概念就化为 \(E \times E\) 上的对称双线性型的概念（A, III, § 6, No.~3）。

我们注意到，条件 (1) 中的第二个与条件 (2) 中的第二个可由其余三个推出。

由 (1) 与 (2) 我们立即得到

\[
f \left(\sum_ {j} \lambda_ {j} x _ {j}, \sum_ {k} \mu_ {k} y _ {k}\right) = \sum_ {j, k} \bar {\lambda} _ {j} \mu_ {k} f \left(x _ {j}, y _ {k}\right).\tag{4}
\]

特别地，若 \(\mathbf{E}\) 是有限维的，且 \((e_j)_{1 \leqslant j \leqslant n}\) 是 \(\mathbf{E}\) 的一个基，则对 \(x = \sum_{j=1}^{n} \xi_j e_j\) 与 \(y = \sum_{j=1}^{n} \eta_j e_j\) ，有

\[
f (x, y) = \sum_ {j, k} \alpha_ {j k} \overline {{{{\xi}}}} _ {j} \eta_ {k}
\]

其中记 \(\alpha_{jk} = f(e_j, e_k)\) ；此外，关系式 (3) 等价于 \(\alpha_{jk} = \overline{\alpha_{kj}}\) 对每对指标 \(j, k\) 成立；这特别蕴含诸数 \(\alpha_{jj}\) 是实数。

由 (3)，数 \(Q(x) = f(x, x)\) 对所有 \(x \in E\) 是实数。此外，我们立即可以建立下列公式，称为极化公式

\begin{enumerate}
\def\labelenumi{(\arabic{enumi})}
\setcounter{enumi}{4}
\tightlist
\item
\end{enumerate}

\[
4 f (x, y) = \sum_ {\varepsilon^ {2} = 1} \varepsilon \mathrm{Q} (x + \varepsilon y) \quad \text { if } \quad \mathbf {K} = \mathbf {R},\tag{6}
\]

\[
4 f (x, y) = \sum_ {\varepsilon^ {4} = 1, \varepsilon \in \mathbf {C}} \varepsilon \mathrm{Q} (x + \overline {{{\varepsilon}}} y) \quad \text { if } \quad \mathbf {K} = \mathbf {C}.
\]

注. --- 我们指出，公式 (6) 对 \(E \times E\) 上的每个半双线性型都成立（也就是说，对每个满足 (1) 与 (2) 但不必满足 (3) 的函数 f 成立）。这个注表明，当 K = C 时，使 \(f(x, x)\) 对所有 \(x \in E\) 为实数的半双线性型 f 必是埃尔米特的：这时关系式 (6) 给出 \(\overline{f(y, x)} = f(x, y)\) ，因为我们有 \(y + \varepsilon x = \varepsilon(x + \overline{\varepsilon}y)\) ，且 \(\mathbf{Q}(\varepsilon z) = \mathbf{Q}(z)\) 对 \(\varepsilon^{4} = 1\) 成立。

由极化公式，我们特别有

命题 1. --- 若 \(f\) 是 \(\mathbf{E}\) 上的埃尔米特型，且 \(\mathbf{M}\) 是 \(\mathbf{E}\) 的一个使 \(f(x, x) = 0\) 对所有 \(x \in \mathbf{M}\) 成立的向量子空间，则也有 \(f(x, y) = 0\) 对每对点 \(x, y\) （\(\mathbf{M}\) 中的）成立。

设 \(f\) 为 \(\mathbf{E}\) 上的埃尔米特型；集 \(\mathbf{N}\) —— 即全体 \(x \in \mathbf{E}\) ，使 \(f(x, y) = 0\) 对所有 \(y \in \mathbf{E}\) 成立 —— 是 \(\mathbf{E}\) 的一个向量子空间。由 (3) 可知，若 \(x_1 \equiv x_2 \pmod{\mathbf{N}}\) 且 \(y_1 \equiv y_2 \pmod{\mathbf{N}}\) ，则有 \(f(x_1, y_1) = f(x_2, y_2)\) ；因此，在商空间 \(\mathrm{E/N}\) 上，我们定义一个半双线性型 \(f\) ，即令 \(\dot{f}(\dot{x}, \dot{y}) = f(x, y)\) （对所有 \(x \in \dot{x}\) 与所有 \(y \in \dot{y}\) ）；显然 \(\dot{f}\) 是埃尔米特的，且关系 \(\langle \dot{f}(\dot{x}, \dot{y}) = 0 \text{ 对所有 } \dot{y} \in \mathrm{E/N} \rangle\) 蕴含 \(\dot{x} = 0 \text{ 于 } \mathrm{E/N}\) ，换言之（A, IX），\(\dot{f}\) 是非退化的。我们说 \(f\) 是与 \(f\) 相伴的非退化埃尔米特型。

